\documentclass[aspectratio=169,14pt]{beamer}
\usepackage[T1]{fontenc}
\usepackage{lmodern,amsmath,amssymb,booktabs,tabularx,tikz}
\usefonttheme{professionalfonts}
\definecolor{ink}{HTML}{24292E}
\definecolor{blue}{HTML}{304E68}
\definecolor{muted}{HTML}{61676B}
\definecolor{repair}{HTML}{A94722}
\setbeamercolor{normal text}{fg=ink,bg=white}
\setbeamercolor{structure}{fg=blue}
\setbeamercolor{frametitle}{fg=ink,bg=white}
\setbeamercolor{alerted text}{fg=blue}
\setbeamerfont{frametitle}{size=\Large,series=\mdseries}
\setbeamertemplate{navigation symbols}{}
\setbeamertemplate{itemize items}[circle]
\setbeamersize{text margin left=9mm,text margin right=9mm}
\setlength{\parskip}{5pt}
\newif\ifbackup
\setbeamertemplate{footline}{\hfill
  {\color{muted}\fontsize{8}{9}\selectfont
  \ifbackup Appendix \number\numexpr\insertframenumber-9\relax
  \else\insertframenumber\fi}\hspace{9mm}\vspace{3mm}}
\newcommand{\DA}{\operatorname{DA}}
\newcommand{\BP}{\operatorname{BP}}
\newcommand{\supp}{C_{\mathrm{support}}}
\newcommand{\resp}{C_{\mathrm{DA}}}
\newcommand{\sources}[1]{\vfill{\fontsize{9}{11}\selectfont\color{muted}#1\par}}
\newcommand{\matchpic}[1]{%
\begin{tikzpicture}[x=1.7cm,y=.64cm,
every node/.style={circle,draw=ink,fill=white,minimum size=4.8mm,inner sep=0pt,font=\small}]
\node (a) at (0,2) {$a$}; \node (b) at (0,1) {$b$}; \node (c) at (0,0) {$c$};
\node (x) at (1,2) {$x$}; \node (y) at (1,1) {$y$}; \node (z) at (1,0) {$z$};
\ifnum#1=2
 \draw[ink,line width=1pt] (a)--(x) (b)--(y) (c)--(z);
 \draw[repair,dashed,line width=1.3pt] (c)--(x);
\else
 \draw[ink,line width=1pt] (a)--(y) (b)--(x) (c)--(z);
\fi
\end{tikzpicture}}
\hypersetup{colorlinks=true,urlcolor=blue,
pdftitle={Housing matches with hidden options},
pdfauthor={Ishan Khillan, Elaine Zhang, Todd Zhou}}
\title{Housing matches with hidden options}
\author{Ishan Khillan, Elaine Zhang, Todd Zhou}
\date{October 29, 2026}
\begin{document}

\begin{frame}[plain]
\vspace{7mm}
{\fontsize{25}{29}\selectfont Housing matches\\with hidden options\par}
\vspace{8mm}
Ishan Khillan, Elaine Zhang, Todd Zhou\par
{\small COMPSCI 2880 \quad October 29, 2026\par}
\vspace{8mm}
\small
Our starting point: brokers can affect which apartments people get to consider.
We want to understand what that does to matching.
\vfill
{\color{muted}Work in progress, for discussion\par}
\end{frame}

\begin{frame}[t]{The broker idea, simplified}
\small
Suppose a renter and a landlord would both prefer to match, but their
agents never put them in contact.

\medskip
We start with renters and listings as the two sides. Each has a fixed
agent who controls access to some of its possible matches.

\medskip
An \emph{audit} opens the feasible matches involving that agent's clients.
It has a cost.

\medskip
This leaves out the reason an agent hides an option.
We have not modeled commissions, friendships, or collusion yet.
\end{frame}

\begin{frame}[t]{What we are assuming}
\small
$E$ contains all feasible renter--listing pairs. Initially, only
$E_0\subseteq E$ is available. Preferences are strict and fixed.

Auditing a set of agents $A$ gives
\[
 E(A)=E_0\cup\{r\ell\in E:f(r)\in A\text{ or }f(\ell)\in A\},
\]
where $f(i)$ is client $i$'s agent. The platform then runs
\[
 D(A)=\DA_R(E(A)).
\]
We want the cheapest $A$ for which no pair in \emph{all of $E$}
would rather match with each other than keep its assignment.

\sources{DA = renter-proposing deferred acceptance. We assume the planner knows
$E$ and the rankings but pays to make pairs available. We still need a housing
setting that justifies these assumptions.}
\end{frame}

\begin{frame}[t]{A small example}
\small
Three renters, three listings. Start with $E_0=\{ay,bx,cz\}$.
Each client has its own agent, and every pair is feasible.
\[
\begin{array}{lll}
 a:x\succ y\succ z & b:y\succ x\succ z & c:x\succ z\succ y\\
 x:b\succ c\succ a & y:a\succ b\succ c & z:c\succ a\succ b
\end{array}
\]
\vspace{-3mm}
\begin{columns}[T,onlytextwidth]
\begin{column}{.31\textwidth}\centering
Start here\par\vspace{1mm}
\matchpic{1}\par
Stable
\end{column}
\begin{column}{.36\textwidth}\centering
Audit $a,b$\par\vspace{1mm}
\matchpic{2}\par
{\color{repair}$cx$ blocks}
\end{column}
\begin{column}{.31\textwidth}\centering
Also audit $c$\par\vspace{1mm}
\matchpic{3}\par
Stable again
\end{column}
\end{columns}
\medskip
In the middle, $c$ wants $x$ instead of $z$, and $x$ wants $c$ instead of $a$.
\sources{The solid lines show DA's assignments. The dashed pair is still hidden.
Preferences have not changed.}
\end{frame}

\begin{frame}[t]{Having a stable matching available isn't enough}
\small
In the middle graph, the original stable matching $\{ay,bx,cz\}$
is still available. DA just doesn't choose it.

\medskip
So there are two different costs:
\begin{align*}
 \supp&:\ \text{make some full-market stable matching available},\\
 \resp&:\ \text{make DA's chosen matching full-market stable}.
\end{align*}
Starting from $N$ copies of that middle graph, plus one hidden isolated pair,
\[
 \supp=1,\qquad \resp=N+1.
\]
\sources{Unit audit costs, one client per agent. The proof is in the appendix.
It is a simple construction. We see it as a reason to study the distinction,
and still need to establish what is new.}
\end{frame}

\begin{frame}[t]{Where this sits in the literature}
\small
\begin{itemize}
\setlength{\itemsep}{8pt}
\item Liu et al.\ (2024) already audit racial steering in LLM
housing recommendations.
\item Matching papers study costly interviews, limited information,
and changes to the available pairs.
\item Trabelsi et al.\ (2025) is especially close: pay to add pairs while
a platform keeps using its own maximum-matching rule.
\end{itemize}
\medskip
Our question is narrower: which agents should we audit so that
DA leaves no blocking pair hidden?
\sources{We need to check whether this exact problem, or an equivalent one,
has already been studied. References and links are in the appendix.}
\end{frame}

\begin{frame}[t]{How we could test the housing part}
\small
Use one fixed catalogue of rental listings. Give an AI recommender
household profiles and ask it to return listing IDs. These recommendations
define the initially visible pairs.

\medskip
Then compare targeted audits with random audits at the same budget,
and with full disclosure.

\medskip
Historical listings could supply rents and apartment features. ACS could
help us generate household profiles. Preferences, agent groups, and audit
costs would still need assumptions.

\sources{This study has not been run. Candidate data: Two Sigma rental listings
and ACS PUMS. Keep catalogue and list length fixed, randomize listing order,
and repeat prompts. The existing experiments use synthetic markets.}
\end{frame}

\begin{frame}[t]{Stability may leave affordable options hidden}
\small
For each modeled renter $r$, let $F_r$ be their affordable feasible listings.
Let $G^+$ be the lower-income renters with at least one such option. Measure
\[
 \mathrm{Access}(A)=\frac{1}{|G^+|}\sum_{r\in G^+}
 \frac{|\{\ell\in F_r:r\ell\in E(A)\}|}{|F_r|}.
\]
Let $C(\tau)$ be the cheapest audit plan that makes DA full-market stable
and gives $\mathrm{Access}(A)\geq\tau$.
\[
 \underbrace{C(\tau)-C(0)}_{\text{extra cost of the access target}}
\]
Does this cost become large when affordable housing is scarce?
Does the recommender change it?
\sources{Assume $|G^+|>0$ and $\tau\in[0,1]$. Report renters with no affordable
option separately. This measures available options, not who ultimately gets a home.}
\end{frame}

\begin{frame}[t]{What we'd like your advice on}
\small
The biggest issue for us is the institution. Is there a housing setting
where a fixed matching rule and costly access make sense?

\bigskip
Would you push this toward a theory result, or toward measuring the
cost of affordable access with an actual recommender?

\bigskip
And which paper should we read before going further?
\end{frame}

\backuptrue

\begin{frame}[t]{Why the costs are 1 and \texorpdfstring{$N+1$}{N+1}}
\small
\setlength{\parskip}{1pt}
Start each copy with
\[
 H=\{ax,ay,az,bx,by,bz,cz\}.
\]
\begin{enumerate}
\setlength{\itemsep}{4pt}
\item $\{ay,bx,cz\}\subseteq H$ is stable in the full market.
But DA returns $\{ax,by,cz\}$, and $cx$ blocks.
\item Without auditing $c$ or $x$, the only new pair we can add is $cy$.
It changes nothing: $c$ prefers $z$, and $z$ keeps $c$.
\item Auditing either $c$ or $x$ exposes $cx$. The rejection chain
returns DA to $\{ay,bx,cz\}$.
\end{enumerate}
So each copy costs zero for support and one for DA. A separate hidden
isolated pair costs one under either objective.
\end{frame}

\begin{frame}[t]{A case where the distinction disappears}
\small
Suppose every listing uses the same strict priority over renters.
Then there is a unique stable matching $M^*$: let renters choose in
priority order.
\[
 M^*\subseteq H\subseteq E
 \quad\Longrightarrow\quad
 \DA_R(H)=M^*.
\]
Each renter's choice from the full market is still available at their turn.
So $\supp=\resp$.

\medskip
We only need to cover the hidden pairs in $M^*$. If agents represent
one side each, this is weighted bipartite vertex cover and we can solve
it by minimum cut.
\sources{This follows from standard matching and covering results.
It is a useful check on which assumptions make our problem difficult.}
\end{frame}

\begin{frame}[t]{References}
\fontsize{9.2}{11}\selectfont
\begin{enumerate}
\setlength{\itemsep}{2pt}
\item Liu, So, Hosoi, D'Ignazio (2024).
\href{https://r-t-lab.org/works/housing-gpt-audit/}{Racial Steering by Large Language Models}.
EAAMO.
\item \href{https://ojs.aaai.org/index.php/ICWSM/article/view/7276}{Asplund et al.\ (2020), Auditing Race and Gender Discrimination in Online Housing Markets}.
ICWSM.
\href{https://doi.org/10.1215/00703370-9357518}{Besbris et al.\ (2021), The Unequal Availability of Rental Housing Information Across Neighborhoods}.
Demography.
\item \href{https://doi.org/10.1016/j.disopt.2016.03.002}{Cseh and Manlove (2016), Stable Marriage and Roommates Problems with Restricted Edges}.
Discrete Optimization.
\href{https://doi.org/10.1613/jair.1.12755}{Boehmer et al.\ (2021), Bribery and Control in Stable Marriage}.
JAIR.
\item \href{https://arxiv.org/abs/1602.04792}{Rastegari, Goldberg, Manlove (2016), Preference Elicitation in Matching Markets via Interviews}.
AAMAS extended abstract and full preprint.
\item \href{https://arxiv.org/abs/2508.18325}{Trabelsi et al.\ (2025), Facilitating Matches on Allocation Platforms}.
ECAI.
\item \href{https://arxiv.org/abs/2609.09693}{Aguiar and Hong (2026), The Attention Cost of Stable Matching}.
Preprint.
\item \href{https://doi.org/10.1007/s10107-003-0375-9}{Hooker and Ottosson (2003), Logic-Based Benders Decomposition}.
\href{https://doi.org/10.1007/BF01586041}{Rothblum (1992), Characterization of Stable Matchings as Extreme Points of a Polytope}.
Mathematical Programming.
\item Data:
\href{https://www.kaggle.com/competitions/two-sigma-connect-rental-listing-inquiries/data}{Two Sigma Connect: Rental Listing Inquiries},
\href{https://www.census.gov/programs-surveys/acs/microdata.html}{ACS Public Use Microdata}.
\end{enumerate}
\end{frame}
\end{document}
